\subsection{Autentifikácia a správa používateľov}
Aplikácia je navrhnutá podľa architektúry \textbf{MVVM (Model-View-ViewModel)}, ktorá zabezpečuje jasné oddelenie biznis logiky od používateľského rozhrania. Implementácia autentifikácie prebieha na úrovni vrstvy \textit{Model} (služba \texttt{AuthService}), zatiaľ čo validácia a spracovanie stavov sú riešené vo vrstve \textit{ViewModel}.

\subsubsection{Implementácia kľúčových funkcií}
Služba \texttt{AuthService} zabezpečuje nízkoúrovňovú komunikáciu s Firebase. Využíva asynchrónne volania pomocou Kotlin korutín, čo umožňuje neblokujúce vykonávanie sieťových operácií. Registrácia nového používateľa zahŕňa nielen vytvorenie účtu, ale aj okamžitú aktualizáciu profilu o zobrazované meno (\textit{displayName}):

\medskip % Pridaná medzera pred kódom
\begin{lstlisting}[language=Kotlin, caption=Registrácia používateľa v AuthService]
suspend fun registerWithEmail(
    email: String,
    password: String,
    displayName: String
): FirebaseUser? {
    val credential = auth.createUserWithEmailAndPassword(email, password).await()
    credential.user?.updateProfile(
        userProfileChangeRequest { this.displayName = displayName }
    )?.await()
    return credential.user
}
\end{lstlisting}
\medskip % Pridaná medzera za kódom

Prihlásenie cez Google účet prebieha prostredníctvom \texttt{GoogleSignInClient}, pričom po získaní \texttt{idToken} sa volá metóda \texttt{firebaseAuthWithGoogle}:

\medskip
\begin{lstlisting}[language=Kotlin, caption=Prihlásenie cez Google v AuthService]
suspend fun firebaseAuthWithGoogle(idToken: String): FirebaseUser? {
    val credential = GoogleAuthProvider.getCredential(idToken, null)
    val result = auth.signInWithCredential(credential).await()
    return result.user
}
\end{lstlisting}
\medskip

Okrem externých poskytovateľov je implementované aj štandardné prihlásenie pomocou e-mailu a hesla:

\medskip
\begin{lstlisting}[language=Kotlin, caption=Prihlásenie pomocou e-mailu a hesla]
suspend fun signInWithEmail(email: String, password: String) {
    auth.signInWithEmailAndPassword(email, password).await()
}
\end{lstlisting}
\medskip

\subsubsection{Validácia vstupných údajov}
Pred odoslaním požiadavky na server vykonáva \texttt{ViewModel} lokálnu validáciu na strane klienta. Týmto prístupom šetríme sieťovú prevádzku a zlepšujeme odozvu rozhrania. V rámci bezpečnosti sa kontroluje:

\begin{itemize}[noitemsep]
    \item formát e-mailovej adresy (validita syntaxe),
    \item komplexnosť hesla (minimálna dĺžka 6 znakov, prítomnosť veľkého písmena a špeciálneho znaku),
    \item zhoda hesla s jeho potvrdením pri registrácii.
\end{itemize}

\medskip
\begin{lstlisting}[language=Kotlin, caption=Validácia hesla v RegisterViewModel]
passwordError = when {
    password.length < 6 -> "Heslo musí mať aspoň 6 znakov."
    !password.any { it.isUpperCase() } -> "Heslo musí obsahovať veľké písmeno."
    !password.any { !it.isLetterOrDigit() } -> "Heslo musí obsahovať špeciálny znak."
    else -> null
}
\end{lstlisting}
\medskip

\subsubsection{Bezpečnostné obmedzenia a banovanie}
Nad rámec štandardnej autentifikácie aplikácia implementuje kontrolu stavu používateľa. Pri každom prihlásení sa overuje atribút \texttt{banned} v dokumente používateľa v databáze Firestore. V prípade, že je používateľ zablokovaný, systém ho okamžite odhlási a zamedzí prístup k chránenému obsahu aplikácie.